% Living Showcase: Einstieg und Navigationsübersicht (Palette Slot 0).
\ZSFTitleHeader{ZSF Template}{\ZSFOwnerDisplayName}
\StartFrontChapter{Living Showcase}

\begin{runintext}
  Diese Ausgabe ist zugleich Präsentation, Bauteilkatalog und visueller
  Regressionstest. Jede Box zeigt sich selbst: Titel und Inhalt beschreiben,
  welcher Baustein hier steht und was er leistet.
\end{runintext}

% Kein \ZSFindex hier: Front-Matter ist unnummeriert und könnte nur auf die
% blosse Kapitelnummer zeigen statt auf einen Abschnitt (rules/55_index).
\SubsectionBar*{Schnellzugriff}

\begin{tablebox}[Welcher Baustein wo\ZSFTitleTag{Übersicht}]
  \begin{ZSFtable}{Y{1.05} Z{0.35} Y{1.60}}
    \ZSFheaderRow \ZSFhead{Bereich} & \ZSFhead{Kap.} & \ZSFhead{Bausteine} \\
    Aussagen & \ZSFsectionref{ch:boxes} & defbox, statementbox, warnbox \\
    Verfahren & \ZSFsectionref{ch:lists} & procedure, factlist \\
    Formeln & \ZSFsectionref{ch:math} & formulabox, goalbox \\
    Tabellen & \ZSFsectionref{ch:tables} & ZSFtable, valuegrid \\
    Visuals & \ZSFsectionref{ch:figures} & ZSFfig, figbox, splitbox \\
    Navigation & \ZSFsectionref{ch:semantics} & Marker, Verweise \\
    Code & \ZSFsectionref{ch:code} & codebox, CodeLine \\
    Stilmittel & \ZSFsectionref{ch:stilmittel} & Farbverfolgung, Klammern
  \end{ZSFtable}
\end{tablebox}

\begin{statementbox}[Lesart]
  Der rechts gesetzte Titel-Tag nennt jeweils die verwendete Umgebung, der
  Titel sagt, wofür sie da ist. So lassen sich Varianten direkt vergleichen
  und auswählen.
\end{statementbox}

\begin{warnbox}[Kein Fachinhalt]
  Die Beispiele tragen bewusst keinen Lernstoff. Sie zeigen ausschliesslich
  das Verhalten der Bausteine — was hier steht, beschreibt den Baustein.
\end{warnbox}
